% SPDX-License-Identifier: CC-BY-SA-4.0
% Author: Matthieu Perrin

\begingroup

\begin{exercice}[Révisions -- Automates finis]\label{exo:lexing/automata/revision}

  On considère l'automate fini $A = \left\langle \{a, b\}, \{0,1,2,3\}, \{0\}, \{2\}, \left\{\begin{array}{lll}
    \langle 0, a, 1\rangle, &
    \langle 1, b, 2\rangle, &
    \langle 2, b, 2\rangle, \\
    \langle 2, a, 3\rangle, &
    \langle 3, a, 3\rangle, &
    \langle 3, b, 3\rangle
  \end{array}\right\} \right\rangle$.

  \begin{question}

  \item Donner la représentation graphique et la représentation matricielle de $A$.

  \item Simuler l'exécution de $A$ sur les mots suivants et dire s'ils sont acceptés :
    $$ab,\quad abb,\quad aab,\quad \varepsilon.$$

  \item L'automate $A$ est-il déterministe ? Est-il complet ? 
    
  \item Décrire en français le langage $\mathcal{L}(A)$.

  \end{question}

  Construire un automate fini déterministe reconnaissant chacun des langages suivants :
  
  \begin{question}

  \item les mots de $\{a,b\}^\star$ se terminant par le facteur $aba$ ;

  \item les mots de $\{a,b\}^\star$ ne contenant pas le facteur $bb$ ;

  \end{question}

\end{exercice}

\begin{correction}

  \begin{question}

  \item 
    \begin{tikzpicture}[automaton, size=20mm, baseline=(0.base)]
      \state[initial]   (0) at (0,0) {$0$};
      \state            (1) at (1,0) {$1$};
      \state[accepting] (2) at (2,0) {$2$};
      \state            (3) at (3,0) {$3$};

      \path (0) edge             node {$a$} (1);
      \path (1) edge             node {$b$} (2);
      \path (2) edge[loop above] node {$b$} (2);
      \path (2) edge             node {$a$} (3);
      \path (3) edge[loop above] node {$a,b$} (3);
    \end{tikzpicture}
    
    $\begin{array}[t]{c|cc}
      &  a & b \\
      \hline
      0& 1 & \\
      1&   & 2\\
      2& 3 & 2\\
      3& 3 & 3
    \end{array}$

  \item 
    \begin{itemize}
    \item $ab$ est accepté  : $\langle ab, 0 \rangle \leadsto \langle b, 1\rangle \leadsto \langle \varepsilon, 2\rangle.$
    \item $abb$ est accepté : $\langle abb, 0 \rangle \leadsto \langle bb, 1\rangle \leadsto \langle b, 2\rangle \leadsto \langle \varepsilon, 2\rangle.$
    \item $aab$ est rejeté : $\langle aab, 0 \rangle \leadsto \langle ab, 1\rangle.$
    \item $\varepsilon$ est rejeté : $\langle \varepsilon, 0 \rangle.$
    \end{itemize}

  \item L'automate $A$ est déterministe mais pas complet.

  \item $\mathcal{L}(A) = ab^+$.

  \item
    \begin{tikzpicture}[automaton, size=20mm, baseline=(0.base)]
      \state[initial]   (0) at (0,0) {$0$};
      \state            (1) at (1,0) {$1$};
      \state            (2) at (2,0) {$2$};
      \state[accepting] (3) at (3,0) {$3$};

      \path (0) edge node {$a$} (1);
      \path (0) edge[loop above] node       {$b$} (0);
      \path (1) edge[loop above] node       {$a$} (1);
      \path (1) edge             node       {$b$} (2);
      \path (2) edge             node       {$a$} (3);
      \path (2) edge[bend left]  node       {$b$} (0);
      \path (3) edge[bend right] node[swap] {$a$} (1);
      \path (3) edge[bend left]  node       {$b$} (2);
    \end{tikzpicture}

  \item
    \begin{tikzpicture}[automaton, size=20mm, baseline=(0.base)]
      \state[initial, accepting] (0) at (0,0) {$0$};
      \state                     (1) at (1,0) {$1$};
      \state                     (2) at (2,0) {$2$};

      \path (0) edge[loop above] node {$a$} (0);
      \path (0) edge[bend left]  node {$b$} (1);
      \path (1) edge[bend left]  node {$a$} (0);
      \path (1) edge             node {$b$} (2);
      \path (2) edge[loop right] node {$a, b$} (2);
    \end{tikzpicture}

  \end{question}
  
\end{correction}

\endgroup
\endinput
